\loai{Loại 1. Sự nóng chảy}
\begin{ex}
	Sự nóng chảy là sự chuyển từ	
	\choice
	{\True thể rắn sang thể lỏng}
	{thể rắn sang thể khí}
	{thể lỏng sang thể rắn}
	{thể lỏng sang thể khí}
	\loigiai{
		\begin{itemize}
			\item Nóng chảy là chuyển từ thể rắn sáng thể lỏng.
		\end{itemize}
	}
	\haidong{2}\end{ex}
\begin{ex}
	Hiện tượng nào \textbf{không} liên quan đến sự nóng chảy?
	\choice
	{Bỏ cục nước đá vào một cốc nước}
	{Đốt một ngọn nến}
	{\True Đốt một ngọn đèn dầu}
	{Đúc một cái chuông đồng}
	\loigiai{
		Đốt một ngọn đèn dầu thì dầu bốc cháy không liên quan tới nóng chảy.
	}
	\haidong{4}\end{ex}
\begin{ex}%[1P7Y1-1][Loai2]
	Chất rắn vô định hình \textbf{không} có đặc tính nào dưới đây?
	\choice
	{Không bị ăn mòn}
	{Dễ định hình}
	{\True Nhiệt độ nóng chảy cao}
	{Không bị gỉ}
	\loigiai{
		\begin{itemize}
			\item Chất rắn vô định hình không có nhiệt độ nóng chảy xác định.
		\end{itemize}
	}%<MyLT1>
	\haidong{2}\end{ex}
\begin{ex}
	Phát biểu nào sau đây \textbf{không} đúng khi nói về sự nóng chảy?
	\choice
	{Sự nóng chảy là sự chuyển từ thể rắn sang thể lỏng của một chất}
	{Mỗi chất khác nhau thường có nhiệt độ nóng chảy khác nhau}
	{Nhiệt độ mà tại đó chất rắn bắt đầu nóng chảy được gọi là nhiệt độ nóng chảy}
	{\True Trong suốt quá trình nóng chảy, nhiệt độ nóng chảy luôn luôn thay đổi}
	\loigiai{
		\begin{itemize}
			\item Trong suốt quá trình nóng chảy, nhiệt độ nóng chảy không thay đổi.
		\end{itemize}
	}
	\haidong{2}\end{ex}
\begin{ex}
	Hiệu ứng nhà kính làm Trái Đất nóng lên, băng ở hai cực tan ra. Băng tan là quá trình nào sau đây?
	\choice
	{\True Quá trình nóng chảy}
	{Quá trình đông đặc}
	{Quá trình sôi}
	{Quá trình bay hơi}
	\loigiai{
		Băng tan là quá trình nóng chảy.
	}
	\haidong{2}\end{ex}
\begin{ex}
	Trong thời gian sắt đông đặc, nhiệt độ của nó
	\choice
	{không ngừng tăng}
	{không ngừng giảm}
	{mới đầu tăng, sau giảm}
	{\True không đổi}
	\loigiai{
		
	}
	\haidong{2}\end{ex}
\loai{Loại 2. Sự đông đặc}
\begin{ex}
	Vào những ngày rất lạnh, nước ở nhiều khu vực của nước ta như Sa Pa, Mẫu Sơn có thể bị đóng băng. Hiện tượng này thể hiện sự chuyển thể nào của chất?
	\choice
	{Sự nóng chảy}
	{\True Sự đông đặc}
	{Sự hóa hơi}
	{Sự ngưng tụ}
	\loigiai{
		\begin{itemize}
			\item Ở những ngày rất lạnh, nhiều khu vực ở nước ta như Sapa, Mẫu Sơn nước có thể bị đóng băng. Hiện tượng này là sự đông đặc.
		\end{itemize}
	}
	\haidong{2}\end{ex}
\loai{Loại 3. Nóng chảy và đông đặc}
\begin{ex}
	Cho bảng theo dõi nhiệt độ của một chất rắn như sau:
	$$
	\begin{array}{|l|c|c|c|c|c|c|}
		\hline \text { Thời gian (phút) } & 0 & 2 & 4 & 6 & 8 & 10 \\
		\hline \text { Nhiệt độ }\left({ }^{\circ} {C}\right) & 20 & 40 & 60 & 80 & 80 & 85 \\
		\hline
	\end{array}
	$$
	Chất rắn nóng chảy trong khoảng thời gian nào?
	\choice
	{\True Từ phút thứ 6 đến phút thứ 8}
	{Từ phút thứ 2 đến phút thứ 4}
	{Từ phút thứ 8 đến phút thứ 10}
	{Từ phút thứ 4 đến phút thứ 10}
	\loigiai{
		
		\begin{itemize}
			\item Chất rắn nỏng chảy trong khoảng thời gian từ phút thứ 6 đến phút thứ 8 (vì nhiệt độ không đổi là $80^\circ C$)
		\end{itemize}
		
	}
	\haidong{3}\end{ex}
\begin{ex}%[1P7B4-1][Loai3]
	Trong sự nóng chảy và đông đặc của các chất rắn kết tinh thì
	\choice
	{mỗi chất rắn nóng chảy ở một nhiệt độ xác định, không phụ thuộc vào áp suất bên ngoài}
	{nhiệt độ đông đặc của chất rắn kết tinh không phụ thuộc vào áp suất bên ngoài}
	{\True mỗi chất rắn kết tinh nóng chảy và đông đặc ở cùng một nhiệt độ xác định trong điều kiện áp suất xác định}
	{mỗi chất rắn kết tinh nóng chảy và đông đặc ở nhiệt độ khác nhau}
	\loigiai{
		\begin{itemize}
			\item Mỗi chất rắn kết tinh nóng chảy và đông đặc ở cùng một nhiệt độ xác định trong điều kiện áp suất xác định.
		\end{itemize}
	}
	\haidong{2}\end{ex}
\begin{ex} \nguon{1}
	Trong quá trình nóng chảy và đông đặc của các chất rắn thì
	\choice
	{mỗi chất rắn nóng chảy ở một nhiệt độ xác định, không phụ thuộc vào áp suất bên ngoài}
	{nhiệt độ đông đặc của chất rắn kết tinh không phụ thuộc vào áp suất bên ngoài}
	{\True với mỗi áp suất xác định, mỗi chất rắn kết tinh nóng chảy và đông đặc ở cùng một nhiệt độ xác định}
	{mỗi chất rắn nóng chảy ở nhiệt độ nào thì cũng sẽ đông đặc ở nhiệt độ đó}
	\loigiai{
		\begin{itemize}
			\item Phát biểu đúng là: \emph{Mỗi chất rắn kết tinh nóng chảy và đông đặc ở cùng một nhiệt độ xác định trong điều kiện áp suất xác định}.
			
			\item Giải thích chi tiết:
			\begin{itemize}
				\item \textbf{Mỗi chất rắn nóng chảy ở một nhiệt độ xác định, không phụ thuộc vào áp suất bên ngoài}: Sai. Nhiệt độ nóng chảy phụ thuộc vào áp suất bên ngoài.
				\item \textbf{Nhiệt độ đông đặc của chất rắn kết tinh không phụ thuộc vào áp suất bên ngoài}: Sai. Nhiệt độ đông đặc cũng phụ thuộc vào áp suất bên ngoài.
				\item \textbf{Mỗi chất rắn kết tinh nóng chảy và đông đặc ở cùng một nhiệt độ xác định trong điều kiện áp suất xác định}: Đúng. Trong điều kiện áp suất xác định, điểm nóng chảy và đông đặc của một chất rắn kết tinh là như nhau.
				\item \textbf{Mỗi chất rắn nóng chảy ở nhiệt độ nào thì cũng sẽ đông đặc ở nhiệt độ đó}: Cần làm rõ điều kiện áp suất xác định, vì phát biểu này có thể lập lờ.
		\end{itemize}\end{itemize}
	}
\end{ex}

\begin{ex} \nguon{2}
	Nhiệt độ nóng chảy của chất rắn kết tinh \textbf{không} có đặc điểm nào sau đây?
	\choice
	{Chất vô định hình không có nhiệt độ nóng chảy xác định}
	{Chất rắn kết tinh có nhiệt độ nóng chảy xác định}
	{\True Thể tích của tất cả các chất đều tăng khi nóng chảy}
	{Với mỗi cấu trúc tinh thể, nhiệt độ nóng chảy phụ thuộc vào áp suất bên ngoài} 
	\loigiai{
		Thể tích của một số chất, như nước đá, giảm khi nóng chảy thay vì tăng.
	}
\end{ex}
\loai{Loại 4. Giải thích hiện tượng nóng chảy - đông đặc}
\begin{ex}
	\immini[thm]{Đồ thị hình vẽ bên biểu diễn sự thay đổi nhiệt độ của nước theo thời gian. Nội dung nào đúng?
		\choice
		{\True Quá trình nóng chảy diễn ra từ phút thứ 1 đến phút thứ 2,5}
		{Quá trình nóng chảy diễn ra trong 1 phút đầu tiên}
		{Từ phút thứ 1 đến phút thứ 2,5 nước ở thể lỏng}
		{Từ phút thứ 2,5 đến phút thứ 3,5 nước bắt đầu sôi}}{\vspace*{.5 cm} \begin{tikzpicture}[scale=1,draw=Mapcolor, very thick, >=stealth']
			\tkzDefPoints{0/0/o, 4/0/x, 0/1.5/y, 0/-1.5/y', 0/-1/a, 1/0/b, 2.5/0/c, 3.5/1/d}
			\draw  [help  lines, xstep=0.5cm,ystep=0.5cm]  (0,-1.5)  grid  (4,1.5);
			\tkzDrawSegments[very thick,Mapcolor,->](o,x y',y)
			\tkzLabelPoint[left](o){$0$}
			\tkzLabelPoint[above](x){t (phút)}
			\tkzLabelPoint[above](y){t ($^\circ C$)}
			\tkzDrawSegments[very thick, blue](a,b b,c c,d)
			\tkzDrawPoints[size=3,fill=white](a,b,c,d,o)
			\tkzLabelPoint[left](a){$-10$}
			\node[left] at (0,1) {$10$};
			\tkzLabelPoint[below](b){$1$}
			\tkzLabelPoint[below](c){$2,5$}
			\node[below] at (3.5,0) {$3,5$};
	\end{tikzpicture}}	
	\loigiai{
		\begin{itemize}
			\item Từ 1 phút tới 2,5 phút nhiệt độ không đổi bằng $0^\circ C$ nên đây là quá trình nóng chảy
		\end{itemize}
	}
	\haidong{5}\end{ex}
\begin{ex} \nguon{2}
	Khi đun nóng kẽm, chúng mềm ra và nóng chảy dần, phát biểu nào sau đây là đúng?
	\choice
	{Trong thời gian nóng chảy, nhiệt độ của kẽm giảm dần}
	{Trong thời gian nóng chảy, nhiệt độ của kẽm lúc tăng lúc giảm}
	{\True Trong thời gian nóng chảy, nhiệt độ của kẽm không đổi}
	{Trong thời gian nóng chảy, nhiệt độ của kẽm tiếp tục tăng}
	\loigiai{
		
		Trong quá trình nóng chảy, nhiệt độ của kẽm không thay đổi cho đến khi toàn bộ lượng kẽm đã biến thành chất lỏng.
	}
\end{ex}
\begin{ex} \nguon{4}
	Các chất rắn vô định hình như thủy tinh, nhựa, sôcôla không có nhiệt độ nóng chảy xác định vì
	\choice
	{chúng không thể nóng chảy khi bị nung nóng}
	{cấu trúc phân tử của chúng được sắp xếp rất chặt chẽ}
	{\True chúng có cấu trúc ngẫu nhiên, không có trật tự như chất rắn kết tinh}
	{khi bị nung nóng, chúng ngay lập tức bay hơi mà không qua trạng thái lỏng}
	\loigiai{
		Các phân tử hoặc các đơn vị cấu trúc của chất rắn vô định hình không được sắp xếp theo một mẫu đặc biệt như trong chất rắn kết tinh. Thay vào đó, chúng tồn tại dưới dạng một cấu trúc ngẫu nhiên.
		Khi nhiệt độ tăng, các phân tử trong chất rắn vô định hình bắt đầu di chuyển và rung động nhanh hơn. Do trong chất rắn vô định hình không có một sự sắp xếp định hình xác định nào nên không xảy ra sự phá vỡ cấu trúc như trong chất rắn kết tinh. Sự chuyển đổi từ trạng thái rắn sang trạng thái lỏng trong chất rắn vô định hình là một quá trình liên tục nên các chất này không có một nhiệt độ nóng chảy cụ thể.
	}
\end{ex}
\begin{ex} \nguon{4}
	Điều gì xảy ra với cấu trúc của sắt khi bị nung đến nhiệt độ nóng chảy $1530^\circ C$?
	\choice
	{Các nguyên tử sắt liên kết chặt chẽ hơn, làm mạng tinh thể trở nên bền vững hơn}
	{\True Các nguyên tử dao động mạnh hơn, làm mạng tinh thể bị phá vỡ và sắt chuyển sang trạng thái lỏng}
	{Các nguyên tử sắt ngừng dao động, giúp sắt giữ nguyên cấu trúc rắn nhưng ở nhiệt độ cao hơn}
	{Mạng tinh thể của sắt bị phá hủy ngay khi bắt đầu gia nhiệt, làm sắt nóng chảy ngay lập tức}
	\loigiai{
		Ở nhiệt độ bình thường, sắt có cấu tạo mạng tinh thể vững chắc, các nguyên tử sắt chỉ dao động quanh vị trí cân bằng cố định. Lúc này sắt tồn tại ở thể rắn.
		Khi được cung cấp nhiệt lượng từ quá trình nung nóng, các nguyên tử tại các nút mạng sẽ dao động mạnh hơn. Ở nhiệt độ nóng chảy, năng lượng dao động của các nguyên tử đã vượt qua năng lượng liên kết giữa chúng ở mạng tinh thể. Kết quả, mạng tinh thể bị phá vỡ nên các nguyên tử sắt bắt đầu chuyển động tự do trong một không gian lớn hơn và trở thành trạng thái lỏng.
	}
\end{ex}

\begin{ex} \nguon{4}
	Phát biểu nào sau đây về đặc tính của chất rắn kết tinh là \textbf{không} đúng?
	\choice
	{Các phân tử liên kết chặt với nhau và sắp xếp theo một trật tự hình học xác định}
	{\True Không có nhiệt độ nóng chảy xác định}
	{Có cấu trúc tinh thể}
	{Có nhiệt độ nóng chảy xác định}
	\loigiai{
		\begin{itemize}
			\item Chất rắn kết tinh có các hạt nguyên tử, phân tử liên kết chặt với nhau và sắp xếp theo trật tự hình học xác định.
			\item Chất rắn kết tinh có cấu trúc tinh thể và có nhiệt độ nóng chảy xác định.
			\item Phát biểu "Không có nhiệt độ nóng chảy xác định" là không đúng đối với chất rắn kết tinh.
		\end{itemize}
	}
\end{ex}

\loai{Loại 5: Bay hơi}
\begin{ex}
	Sự hóa hơi xảy ra trên bề mặt chất lỏng gọi là
	\choice
	{sự nóng chảy}
	{sự sôi}
	{sự đông đặc}
	{\True sự bay hơi}
	\loigiai{
		Sự hóa hơi xảy ra trên bề mặt chất lỏng gọi là sự bay hơi}
	\haidong{2}\end{ex}
\begin{ex} \nguon{1}
	Khi đang chuyển thể thì quá trình nào xảy ra ở nhiệt độ \textbf{không} xác định?
	\choice
	{Sự nóng chảy}
	{Sự đông đặc}
	{Sự sôi}
	{\True Sự bay hơi}
	\loigiai{
		\begin{itemize}
			\item Phát biểu đúng là: \emph{Sự bay hơi}.		
			\item Giải thích chi tiết:
			\begin{itemize}
				\item \textbf{Sự nóng chảy}: Sai. Nóng chảy của chất rắn kết tinh xảy ra ở nhiệt độ xác định.
				\item \textbf{Sự đông đặc}: Sai. Đông đặc của chất rắn kết tinh cũng xảy ra ở nhiệt độ xác định.
				\item \textbf{Sự sôi}: Sai. Sự sôi xảy ra ở nhiệt độ xác định gọi là nhiệt độ sôi.
				\item \textbf{Sự bay hơi}: Đúng. Bay hơi có thể xảy ra ở bất kì nhiệt độ nào dưới điểm sôi, không có nhiệt độ cố định.
		\end{itemize}\end{itemize}
	}
\end{ex}
\begin{ex}
	Trong các trường hợp dưới đây, trường hợp nào liên quan đến sự bay hơi?
	\choice
	{Kính cửa sổ bị mờ đi trong những ngày đông giá lạnh}
	{\True Cốc nước bị cạn dần khi để ngoài trời nắng}
	{Miếng bơ để bên ngoài tủ lạnh sau một thời gian bị chảy lỏng}
	{Đưa nước vào trong tủ lạnh để làm đá}
	\loigiai{
		\begin{itemize}
			\item 	Cốc nước bị cạn dần khi để ngoài trời nắng liên quan tới sự bay hơi.
		\end{itemize}
	}
	\haidong{3}\end{ex}
\loai{Loại 6: Tốc độ bay hơi}
\begin{ex}
	Khi làm muối từ nước biển, người dân dẫn nước biển vào các ruộng muối rồi để nước biển bay hơi thì sẽ thu được muối. Theo em, thời tiết như thế nào thì thuận lợi cho nghề làm muối?
	
	\choice
	{\True Trời nắng nóng và nhiều gió}
	{Trời ít gió}
	{Không khí ẩm}
	{Trời lạnh}
	\loigiai{
		\begin{itemize}
			\item Trời càng nắng và gió càng nhiều thì nước bay hơi càng nhanh nên nhanh thu được muối hơn
		\end{itemize}
	}
	\haidong{3}\end{ex}
\begin{ex} \nguon{1}
	Phơi quần áo ở nơi có nắng và gió thì quần áo khô nhanh hơn vì
	\choice
	{ánh nắng mặt trời làm nước trong quần áo ngưng tụ nhanh hơn}
	{\True nhiệt độ cao làm tăng tốc độ bay hơi và gió giúp loại bỏ hơi nước nhanh chóng}
	{gió làm tăng nhiệt độ bề mặt quần áo, giúp nước bay hơi nhanh hơn}
	{ánh nắng mặt trời làm giảm áp suất không khí xung quanh, khiến nước bay hơi nhanh hơn}
	\loigiai{
		Vì có hai yếu tố quan trọng ảnh hưởng đến tốc độ khô của quần áo: nhiệt độ và gió.
		\begin{itemize}
			\item Nhiệt độ: Ánh nắng mặt trời tạo ra nhiệt độ cao, làm tăng tốc độ bay hơi của nước trong quần áo. Khi nhiệt độ tăng lên, phân tử nước trong quần áo sẽ di chuyển nhanh hơn và dễ bay hơi, giúp quần áo khô nhanh hơn.
			\item Gió: Gió thổi qua quần áo giúp loại bỏ hơi ẩm bay hơi nhanh chóng. Khi gió thổi qua bề mặt quần áo, nó góp phần tăng lực kéo phân tử chất lỏng từ quần áo ra bên ngoài.
		\end{itemize}
	}
\end{ex}

\loai{Loại 7: Sôi}
\begin{ex}
	Sự sôi là
	
	\choice
	{sự hóa hơi của các chất ở mọi nhiệt độ}
	{sự hóa hơi của các chất xảy ra trên bề mặt chất lỏng ở mọi nhiệt độ}
	{\True sự hóa hơi của các chất xảy ra cả trên bề mặt và trong lòng khối chất lỏng ở một nhiệt độ xác định}
	{sự hóa hơi của các chất xảy ra cả trên bề mặt và trong lòng khối chất lỏng ở mọi nhiệt độ}
	\loigiai{
		\begin{itemize}
			\item Sự sôi của chất là sự hóa hơi của chất xảy ra cả trên bề mặt và trong lòng khối chất lỏng ở một nhiệt độ xác định.
		\end{itemize}
	}
	\haidong{1}\end{ex}
\loai{Loại 8: Nhiệt độ sôi}
\begin{ex}
	Phát biểu nào sau đây là \textbf{sai}?
	\choice
	{Nhiệt độ sôi của chất lỏng phụ thuộc vào áp suất khí phía trên bề mặt chất lỏng}
	{Áp suất khí càng cao thì nhiệt độ sôi của chất lỏng càng cao}
	{\True Áp suất khí càng nhỏ thì nhiệt độ sôi của chất lỏng càng cao}
	{Ở một áp suất nhất định, mỗi chất lỏng sôi ở nhiệt độ xác định và không đổi}
	\loigiai{
		\begin{itemize}
			\item Áp suất khí càng nhỏ thì nhiệt độ sôi của chất lỏng càng thấp.
		\end{itemize}
	}
	\haidong{2}\end{ex}
\begin{ex} \nguon{1}
	Thức ăn trong nồi áp suất thường nhanh mềm hơn so với nấu trong nồi thường vì
	\choice
	{nồi áp suất làm giảm nhiệt độ sôi của nước, giúp thực phẩm nhanh chín}
	{\True nồi áp suất làm tăng áp suất trong nồi, khiến nước sôi ở nhiệt độ cao hơn, giúp thức ăn nhanh mềm hơn}
	{nồi áp suất truyền nhiệt tốt hơn các loại nồi thông thường}
	{trong nồi áp suất, hơi nước ngưng tụ thành nước nhanh hơn, giúp thức ăn chín nhanh}
	\loigiai{
		\begin{itemize}
			\item Ta biết rằng chất lỏng ở các áp suất khác nhau sự sôi diễn ra ở các nhiệt độ khác nhau. 
			\item Vì vậy, muốn cho nước trong nồi sôi ở nhiệt độ sôi cao hơn, nấu chín thức ăn nhanh hơn thì tất phải tăng cao áp suất trong nồi, vì thế ta dùng nồi áp suất.\end{itemize}
	}
\end{ex}
\loai{Loại 9: Hiện tượng thực tế về sôi}
\begin{ex}
	Bình nước sôi muốn để nguội nhanh thì cần mở nắp để hơi nước thoát ra vì
	\choice
	{\True mở nắp giúp hơi nước nóng thoát ra ngoài, mang theo nhiệt lượng và làm nước trong bình nguội nhanh hơn}
	{mở nắp giúp không khí bên ngoài tràn vào, làm nước bốc hơi nhanh hơn}
	{mở nắp làm giảm áp suất bên trong bình, khiến nhiệt độ nước giảm nhanh chóng}
	{mở nắp giúp giảm nhiệt độ sôi của nước, khiến nước mất nhiệt nhanh hơn}
	\loigiai{
		\begin{itemize}
			\item Vì hơi nước có nhiệt độ cao, khi mở nắp thì hơi nước thoát ra nhiều và nhanh hơn làm cho nước còn lại trong bình nguội nhanh hơn.
		\end{itemize}
		
	}
	\haidong{6}\end{ex}
\begin{ex} \nguon{4}
	Quá trình chuyển từ thể khí sang thể lỏng của các chất được gọi là
	\choice
	{sự đông đặc}
	{sự hóa hơi}
	{sự nóng chảy}
	{\True sự ngưng tụ}
	\loigiai{
		Ngưng tụ là quá trình chuyển từ thể khí sang thể lỏng khi nhiệt độ giảm xuống.
	}
\end{ex}
\loai{Loại 10: Hiện tượng ngưng tụ}
\begin{ex}
	Đưa cốc nước lạnh ra ngoài trời nóng thì thấy xuất hiện một lớp nước bám ngoài thành cốc. Đó là do hiện tượng
	\choice
	{bay hơi}
	{nóng chảy}
	{thăng hoa}
	{\True ngưng tụ}
	\loigiai{
		\begin{itemize}
			\item Khi đưa cốc nước lạnh ra ngoài trời nóng, hơi nước trong không khí gặp bề mặt lạnh của cốc sẽ bị ngưng tụ và chuyển thành các giọt nước bám trên thành cốc. Đây là hiện tượng ngưng tụ.
			\item Bay hơi là quá trình ngược lại của ngưng tụ, xảy ra khi chất lỏng chuyển thành hơi.
			\item Nóng chảy là quá trình chuyển từ thể rắn sang thể lỏng.
			\item Thăng hoa là hiện tượng chất rắn chuyển trực tiếp thành hơi mà không qua trạng thái lỏng.
		\end{itemize}
	}
\end{ex}
\begin{ex}
	Hiện tượng tự nhiên nào sau đây là do hơi nước ngưng tụ?
	
	\choice
	{Gió thổi}
	{Lốc xoáy}
	{\True Mưa rơi}
	{Mây bay}
	\loigiai{
		\begin{itemize}
			\item Mưa rơi là quá trình ngưng tụ của hơi nước
		\end{itemize}
	}
	\haidong{3}\end{ex}
\loai{Loại 11: Bay hơi và ngưng tụ}

\loai{Loại 12: Tổng hợp sự chuyển thể}
\begin{ex} \nguon{4}
	Quá trình chuyển từ thể rắn sang thể khí của các chất được gọi là
	\choice
	{sự đông đặc}
	{sự hóa hơi}
	{sự nóng chảy}
	{\True sự thăng hoa}
	\loigiai{
		Sự thăng hoa là quá trình chuyển trực tiếp từ thể rắn sang thể khí mà không qua giai đoạn lỏng.
	}
\end{ex}
\begin{ex} \nguon{4}
	Quá trình chuyển từ thể khí sang thể rắn của các chất được gọi là
	\choice
	{\True sự ngưng kết}
	{sự thăng hoa}
	{sự đông đặc}
	{sự ngưng tụ}
	\loigiai{
		Sự ngưng kết là quá trình chuyển từ thể khí sang thể rắn mà không qua giai đoạn lỏng.
	}
\end{ex}
\loai{Loại 13: Giải thích hiện tượng hóa hơi – ngưng tụ}
\begin{ex}
	Quần áo khô sau khi phơi dưới ánh nắng mặt trời. Hiện tượng này thể hiện?
	\choice
	{\True Sự bay hơi của nước}
	{Sự đông đặc của nước}
	{Sự ngưng tụ của nước}
	{Sự nóng chảy của nước}
	\loigiai{
		
		\begin{itemize}
			\item Quần áo khô sau khi phơi dưới ánh nắng mặt trời thể hiện sự bay hơi của nước.
		\end{itemize}
		
	}
	\haidong{2}\end{ex}
\loai{Loại 14: NHH – Đồ thị thí nghiệm}
\begin{ex} \nguon{2}
	\immini[thm]{Trên hình vẽ là các đồ thị biểu diễn sự thay đổi nhiệt độ theo thời gian của bốn chất. Ban đầu, tất cả các chất này đều ở trạng thái lỏng. Chất nào trong số các chất này có nhiệt độ sôi cao nhất?
		\choice[2]
		{\True 1}
		{2}
		{3}
		{4}}{\begin{tikzpicture}[scale=0.6,transform shape,draw=Mapcolor, very thick, >=stealth']
			\tkzDefPoints{0/0/o, 4/0/x, 0/4/y}
			\tkzDrawSegments[thick,->](o,x o,y)
			\tkzLabelPoint[below left](o){$0$}
			\tkzLabelPoint[below](x){$\tau\ (s)$}
			\tkzLabelPoint[above](y){$t\ (^\circ C)$}
			\draw[very thick] (0,0.5) -- (0.75,1) -- (3,1) -- (3.25,2);
			\draw[very thick] (0,0.5) -- (0.75,1.75) -- (2,1.75) -- (3,2.5);
			\draw [very thick](0,0.5) -- (0.5,2.25) -- (1.25,2.25) -- (1.5,2.75);
			\draw [very thick](0,0.5) -- (0.25,2.75) -- (0.75,2.75) -- (1.25,3.25);
			\node at (1.5,3.5) {1};
			\node at (1.75,3) {2};
			\node at (3.25,2.75) {3};
			\node at (3.5,2) {4};
	\end{tikzpicture}}
	
	\loigiai{
		Chất 1 có đồ thị kéo dài nhất trong trạng thái lỏng trước khi nhiệt độ đạt đến điểm sôi, vì vậy nó có nhiệt độ sôi cao nhất.
	}
\end{ex}
\loai{Loại 15: Câu hỏi đúng sai}
\begin{ex} \nguon{1}
	Trước khi trời nổi cơn giông hoặc sắp mưa thì không khí thường rất oi bức.
	\choiceTF[t]
	{\True Vì quá trình mây ngưng tụ thành mưa tỏa ra nhiệt lượng vào không khí}
	{Vì mây thể hơi sẽ hấp thụ nhiệt lượng từ môi trường xung quanh}
	{Trước khi mưa, không khí trở nên oi bức vì áp suất không khí tăng lên do quá trình ngưng tụ của mây}
	{Trước khi trời nổi cơn giông hoặc sắp mưa, không khí thường rất oi bức do áp suất không khí giảm đi nhanh chóng}
	\loigiai{
		\begin{itemchoice}
			\itemch Vì khi đó đang diễn ra quá trình mây (thể hơi) ngưng tụ thành mưa (thể lỏng) thì nó tỏa một nhiệt lượng vào không khí nên ta cảm thấy nóng nực, oi bức.
			\itemch 
			\itemch 
			\itemch 
		\end{itemchoice}
	}
\end{ex}
\begin{ex} \nguon{1}
	Khi tắm xong ta lại có cảm giác mát lạnh và vừa tắm xong thì không nên đứng trước gió.
	\choiceTF[t]
	{\True Sau khi tắm, nước bám trên da sẽ bay hơi, làm cho cơ thể cảm thấy mát lạnh}
	{\True Quá trình bay hơi lấy nhiệt lượng từ cơ thể, làm giảm nhiệt độ của cơ thể}
	{Nếu đứng trước gió sau khi tắm, quá trình bay hơi trên da diễn ra chậm hơn, làm cho cơ thể giữ nhiệt lượng và cảm thấy ấm áp hơn}
	{\True Đứng trước gió khi vừa tắm xong làm cho quá trình bay hơi trên da diễn ra nhanh hơn, làm mất nhiệt lượng của cơ thể nhanh chóng, dễ dẫn đến cảm lạnh}
	\loigiai{
		\begin{itemchoice}
			\itemch Vì sau khi tắm có nước bám trên da và nước có xu hướng bay hơi. Khi nước bay hơi (chuyển từ thể lỏng sang thể hơi) nó sẽ lấy một phần nhiệt lượng của cơ thể làm nhiệt độ cơ thể giảm xuống một chút gây cho ta cảm giác mát lạnh.
			\itemch 
			\itemch Nếu khi vừa tắm xong mà đứng trước gió làm cho quá trình bay hơi trên da diễn ra nhanh hơn thì sự mất nhiệt của cơ thể sẽ nhanh và nhiều nên ta dễ bị cảm lạnh.
			\itemch 
		\end{itemchoice}
	}
\end{ex}
\begin{ex} \nguon{4}
	\immini[thm]{Những người dân ở những vùng khô nóng (sa mạc) thường lấy khăn ướt quấn quanh cổ chai lọ đựng nước để khi khăn khô họ có nước mát trong chai để uống. Nguyên nhân chủ yếu là do
		\choiceTF[t]
		{tránh được ánh nắng mặt trời}
		{\True quá trình bay hơi của nước trong khăn sẽ lấy đi nhiều nhiệt từ chai nước, làm cho nước trong chai lạnh đi khá nhanh}
		{khăn ướt đã được nhúng qua nước đá}
		{nước trong khăn thu nhiệt từ chai nước để ấm lên}}{
\includegraphics[width=4cm]{Images/khanuot}}
	\loigiai{
		\begin{itemchoice}
			\itemch 
			\itemch 
			\itemch 
			\itemch 
		\end{itemchoice}
		Khi khăn ướt được quấn quanh một chai lọ hoặc bất kì vật thể nào, nước từ khăn bắt đầu bay hơi. Quá trình bay hơi này đòi hỏi cần năng lượng, và năng lượng này được cung cấp từ môi truờng xung quanh. Khi nước bay hơi từ khăn, nó lấy đi một lượng nhiệt từ vật thể, làm giảm nhiệt độ của vật thể đó. Do đó, nhiệt độ của chai lọ và nước bên trong giảm đi, tạo ra cảm giác mát mẻ hơn khi uống.
	}
\end{ex}
\begin{ex} \nguon{3}
	\immini[thm]{Hình bên dưới là đường biểu diễn sự thay đổi nhiệt độ theo thời gian của thí nghiệm đun nóng liên tục của một lượng nước đá trong một bình không kín.}{\begin{tikzpicture}[scale=.6,thick,>=stealth']
			\tkzDefPoints{0/0/o, 6.75/0/x, 0/5/y}
			\tkzDrawSegments[very thick,Mapcolor,->](o,x o,y)
			\tkzLabelPoint[below left](o){$O$}
			\tkzLabelPoint[below](x){$t$ (phút)}
			\tkzLabelPoint[above](y){$t$ ($^\circ C$)}
			\draw[dashed,thin] (3.75,0) -- (3.75,4);
			\draw[dashed,thin] (5.25,0) -- (5.25,4);
			
			\draw[red, very thick] (0,0) -- (1.25,0) -- (3.75,4) -- (5.25,4) -- (6.5,4.5);
			
			\foreach  \x  in  {1,2,...,4}
			{\pgfmathtruncatemacro{\label}{75*\x};
				\draw[fill=white](1.25*\x,0)circle(0.5pt);}
			
			\foreach  \y  in  {1,2,...,4}
			{\pgfmathtruncatemacro{\label}{5*\y};
				\draw[fill=white](0,\y)circle(0.5pt);}
			
			\node[left] at (0,2) {$50$};
			\node[left] at (0,4) {$100$};
			\node[below] at (1.25,0) {$A$};
			
			\node[above] at (3.75,4) {$B$};
			\node[above] at (5.25,4) {$C$};
			\node[above] at (6.5,4.5) {$D$};
	\end{tikzpicture}}
	\choiceTF[t]
	{\True Đoạn $OA$ cho biết nước tồn tại ở cả thể rắn và thể lỏng}
	{\True Đoạn $CD$ cho biết nước không tồn tại ở thể lỏng}
	{Đoạn $AB$ cho biết nước đang tồn tại ở thể rắn}
	{\True Đoạn $BC$ cho biết nước đang sôi}
	\loigiai{
		\begin{itemchoice}
			\itemch Đoạn $OA$ là quá trình nước đá nóng chảy, chuyển từ thể rắn sang thể lỏng, nên giai đoạn này tồn tại cả thể rắn và thể lỏng.
			\itemch Đoạn $AB$ là đoạn nước đá đã nóng chảy hoàn toàn, chuyển thành thể lỏng và đang nóng lên.
			\itemch Đoạn $BC$ là nước đang sôi, và đang chuyển từ thể lỏng sang thể khí.
			\itemch Đoạn $CD$ là nước đã hóa hơi hoàn toàn, và nhiệt độ khí đang tăng lên.
		\end{itemchoice}
	}
\end{ex}
\begin{ex} \nguon{4}
	\immini[thm]{Nước có vai trò rất quan trọng trong cuộc sống. Sự biến đổi các trạng thái của nước tạo nên một vòng tuần hoàn được gọi là vòng tuần hoàn của nước. Hình bên mô tả một cách đơn giản vòng tuần hoàn của nước.}{\includegraphics[width=4cm]{images/2024_05_29_5f55f4890f9e0910bb1cg-14(2)}}
	\choiceTF[t]
	{\True Trong vòng tuần hoàn của nước như hình bên, quá trình $A$ là sự bay hơi của các phân tử nước ở bề mặt sông, hồ,..}
	{\True Trong quá trình $B$, nước đã chuyển từ thể khí sang thể lỏng}
	{Trong sự chuyển thể thành mưa, hơi nước đã phải hấp thu một lượng nhiệt lớn từ không khí}
	{\True Năng lượng cung cấp cho nước thực hiện quá trình bay hơi chủ yếu được cung cấp từ Mặt Trời}
	
	\loigiai{
		\begin{itemchoice}
			\itemch 
			\itemch 
			\itemch Hơi nước đã tỏa ra một lượng nhiệt lớn vào môi trường
			\itemch 
		\end{itemchoice}
	}
\end{ex}
\begin{ex} \nguon{4}
	\immini[thm]{Để tạo ra các loại rượu truyền thống đặc trưng của Việt Nam, các cơ sở sản suất rượu đã thực hiện nhiều công đoạn. Trong các công đoạn đó có một công đoạn rất quan trọng là quá trình chưng cất rượu. Quá trình chưng cất được thể hiện đơn giản bằng sơ đồ như hình bên.}{\includegraphics[width=5cm]{images/2024_05_29_5f55f4890f9e0910bb1cg-14(1)}}
	\choiceTF[t]
	{\True Ở bồn A hỗn hợp nguyên liệu lỏng được cung cấp nhiệt lượng để tạo hơi rượu}
	{\True Ở bồn $B$ diễn ra quá trình hơi rượu ngưng tụ thành rượu dưới dạng lỏng}
	{Nước làm mát ở bồn $B$ có tác dụng cung cấp nhiệt lượng cho quá trình ngưng tụ của hơi rượu}
	{\True Thực ra trong ống dẫn hơi rượu vừa có hơi rượu vừa có hơi nước}
	
	\loigiai{
		\begin{itemchoice}
			\itemch 
			\itemch 
			\itemch Nước mát ở bồn B có tác dụng nhận nhiệt từ hơi rượu để hơi rượu ngưng tụ
			\itemch Ở bồn B khi được cung cấp nhiệt lượng nước cũng có sự hóa hơi cùng với rượu nhưng ít hơn.
		\end{itemchoice}
	}
\end{ex}